Regression-based estimators of the global average treatment effect (GATE) under network interference fit an outcome model on the experimental data and predict ``all treated'' versus ``all control''. How accurate are they, what determines their error, and when do they fail? For least-squares estimators we show an exact decomposition: the estimator is a weighted sum $\sum_i w_iy_i$ whose weights depend only on the realised assignment, so the noise error is exactly $\sigma\|w\|_2$ and the misspecification error is controlled by $\|w\|_1$ and the approximation residual. For Bernoulli designs we derive that the noise error grows like $\sqrt{d/n}$ in the degree $d$, and that the large-degree bias equals the local slope of the exposure--response curve at the treatment probability rather than its secant slope. In 11 simulation experiments (eight interference mechanisms; synthetic graphs, ego-Facebook and Facebook100-Stanford3; Bernoulli, cluster and saturation designs; eleven estimators) the noise prediction holds to within a few percent (ratio $1.00$) and the bias prediction matches the observed bias within Monte-Carlo error in 30 settings, without fitted parameters. The bias does not shrink with $n$; it changes sign with the treatment probability. Cluster randomisation removes most of the extrapolation error only when the graph has small-cut clusters, and then ignoring interference is nearly as good; flexible learners (gradient boosting) and richer feature sets were worse than a two-feature linear regression; regression can be worse than the naive difference in means at high degree or weak interference; edge errors in the analyst's graph bias Bernoulli-design estimates roughly proportionally; and i.i.d./robust standard errors under-cover (estimated SEs as small as 0.32 of the true SD) whereas cluster-robust errors help only when the clusters are meaningful. All conclusions are conditional on our synthetic mechanisms, and we report failures and limitations.
